%% The algebraic rescaling (D. Bothe, 2026-10): order-of-accuracy assumptions, separately for the
%% unstructured finite-volume (PDE) solvers and for the semi-Lagrangian method.
\section{The algebraic rescaling of $\psi$: what accuracy to expect}
\label{sec:rescaling}

\paragraph{The two maps (D.~Bothe).} Both multiply $\psi$ by a positive factor, so the sign of
$\psi$, the zero set $\Sigma$ and the normal on $\Sigma$ stay exactly as they are:
\begin{align}
  \text{(0a)}\quad r &= \frac{\psi}{|\nabla\psi|}, \label{eq:r}\\
  \text{(0b)}\quad \hat d &= \frac{2\psi}{|\nabla\psi| +
     \sqrt{\left(|\nabla\psi|^2 - 2\psi\,\partial_{nn}\psi\right)^+}},
  \qquad \partial_{nn}\psi = \bn\cdot(\nabla\nabla\psi)\,\bn,\ \bn=\frac{\nabla\psi}{|\nabla\psi|}.
  \label{eq:dhat}
\end{align}
(0b) is the stable root of $\psi-d\,\partial_n\psi+\tfrac{d^2}{2}\partial_{nn}\psi=0$, the
Taylor expansion of $\psi(\bx-d\bn)=0$ to second order. With $A=|\nabla\psi|$ and
$B=\partial_{nn}\psi$ on $\Sigma$ (\Der):
\begin{equation}
  r = d - \frac{B}{2A}\,d^2 + O(d^3), \qquad \hat d = d + O(d^3).
\end{equation}
So (0a) is a first-order distance (exact value and slope on $\Sigma$) and (0b) a second-order
distance. Both return $\psi$ unchanged if $\psi=d$. For $\psi=r^2-R^2$,
$|\nabla\psi|^2-2\psi\,\partial_{nn}\psi=4R^2$ and (0b) gives $r-R$ exactly (\Der).

\paragraph{How the error is counted.} In the band $|d|\le kh$ (the sign-change cells and $k-1$
layers) the pointwise error of a discrete map has a model part and a part from the discrete
derivatives. To leading order, with the errors $\delta q$ of $|\nabla\psi|$ and $\delta B$ of
$\partial_{nn}\psi$ (\Der):
\begin{equation}
  \delta r \approx -\,d\,\delta q, \qquad
  \delta\hat d \approx -\,d\,\delta q + \tfrac12 d^2\,\delta B .
  \label{eq:derr}
\end{equation}
The error of the field near $\Sigma$ is reported in units of $h$ (the static test's
\texttt{dist3}), so a pointwise $O(h^{p+1})$ is ``order $p$''.

\subsection{Unstructured finite volumes (the PDE solvers)}
\label{sec:rescaling-fvm}

The PDE solvers have no reconstruction. The UFVM version takes $\nabla\psi$ from a FVM gradient
(least squares or Gauss) and the Hessian as the FVM gradient of that gradient,
$\nabla\nabla\psi\approx\nabla_h(\nabla_h\psi)$ (\path{rescalingRedistancer}, the redistancer
family of \texttt{leiaLevelSetFoam}).

\begin{enumerate}
\item Uniform hexahedra, smooth $\psi$ (\Der). The centred gradient has $\delta q=O(h^2)$, and
the gradient of the gradient is a consistent second difference over $2h$ with $\delta B=O(h^2)$.
By \eqref{eq:derr}: (0a) error $O(h^2)$ from the model term, order~1 in units of $h$; (0b)
error $O(h^3)$, order~2.
\item Noise and the grid-scale mode (\Der, corrected 2026-10-29). A perturbation $\varepsilon$
of $\psi$ gives $\delta q=O(\varepsilon/h)$ and $\delta B=O(\varepsilon/h^2)$; in the band,
\eqref{eq:derr} turns this into $O(k\varepsilon)$ and $O(k^2\varepsilon)$. Linearised about $d$,
the first-order map with the centred gradient is $\delta\leftarrow\delta-d\,(\bn\cdot\nabla_h\delta)$:
exactly one step of the \texttt{linearQ} source with $\mu\Delta t=1$. For the mode
$\delta=E(-1)^{i+j}d$ the centred difference gives $\bn\cdot\nabla_h\delta=-(-1)^{i+j}E$, so one
application DOUBLES the mode. (A first version of this report said the maps are neutral on the
unstable mode; that holds only for the pure staircase, not for the $d$-modulated mode.) Applied
once, the map is accurate; applied repeatedly with the centred gradient, it amplifies the mode by
up to 2 per application (\Meas: applied every step on $d$, $N=64$, it diverges, with
Crank--Nicolson and with Euler). A one-sided (Rouy--Tourin) $|\nabla\psi|$ does not vanish at a
grid-scale extremum and damps the mode instead (\Der{} in the first-campaign report,
sec.~eigenmode); \texttt{magGrad godunov} selects it.
\item Time integration (\Meas, 2026-10-29). A change of $\psi$ outside the linear solve (any
redistancer event) is re-applied as a rate by Crank--Nicolson, which stores the last increment:
one rescaling of $r^2-R^2$ at $N=128$ kept half of its error after 128 further steps with
Crank--Nicolson, and reached round-off ($9.7\times10^{-15}h$) with Euler. The rescaling tests
therefore integrate with Euler (\texttt{PSI\_DDT\_SCHEME}).
\item Non-uniform and polyhedral meshes (\Der, to be measured). The FVM gradient is first order,
$\delta q=O(h)$, and the gradient of the gradient is not consistent, $\delta B=O(1)$. Then
$\delta\hat d=O(h^2)$, as $\delta r$: (0b) loses its advantage and both maps are order~1. This
is the prediction the UFVM Hessian is built to test (author decision 2026-10-29: implement it
also to show whether its noise destroys the solution).
\item Repeated application. The discrete map is not a projection: a second application changes
the field again by the discretisation error. Applied every step for $\tau/\Delta\tau=O(N)$
steps, a per-event error $O(h^{p+1})$ can accumulate to $O(h^p)$ unless the map contracts. The
static test measures one event and one event per step (pre-registered in the configs).
\end{enumerate}

\subsection{The semi-Lagrangian method (quadratic fit)}
\label{sec:rescaling-sl}

The semi-Lagrangian method fits a weighted least-squares quadratic per cell. At the cell centre
the fit gives $\nabla\psi$ and $\nabla\nabla\psi$ from one compact stencil.

\begin{enumerate}
\item Any mesh, smooth $\psi$, a well-conditioned fit (\Der). A consistent quadratic fit has
$\delta q=O(h^2)$ and $\delta B=O(h)$, also on polyhedra. By \eqref{eq:derr}: (0a) order~1 (the
model term), (0b) $O(h^3)$, order~2, on every mesh. So the reconstruction keeps the advantage of
(0b) where the UFVM Hessian loses it.
\item Noise (\Der). The same scaling as in the UFVM case, multiplied by the fit's amplification
factor $\Lambda$: the amplification report measured $\Lambda_{\max}=1.05$ on blockMesh
hexahedra and $1.26$ on pMesh polyhedra. The fit is not blind to the odd-even mode in the same
way as the centred difference; how much of the mode the fit removes is to be measured.
\item The value. The map multiplies the cell value $\psi_c$, not the fit value at the centre
(the least-squares fit does not interpolate), so the sign of $\psi_c$ is kept exactly.
\end{enumerate}
